% Figura: líneas de campo y superficies equipotenciales (corte meridiano).
\begin{figure}[ht]
\centering
\IfFileExists{rapido.flag}{\fbox{\parbox{0.9\textwidth}{\centering\small Figura omitida (modo rápido)}}}{%
\newcommand{\marcoL}{\draw[grispizarra, line width=0.5pt] (-3,-3) rectangle (3,3);}
\tikzset{lcampo/.style={cm1, line width=0.55pt, postaction={decorate}, decoration={markings, mark=at position 0.5 with {\arrow{Stealth[length=3.6pt]}}}},
         lequi/.style={azulacero, line width=0.5pt}}
\begin{minipage}[t]{0.32\textwidth}\centering
\begin{tikzpicture}[x=0.8cm, y=0.8cm, baseline=(current bounding box.north)]
  \begin{scope}
    \clip (-3,-3) rectangle (3,3);
    \foreach \r in {0.5,1.0,1.5,2.0,2.5,3.0}{\draw[lequi] (0,0) circle (\r);}
    \foreach \a in {0,22.5,...,337.5}{%
      \draw[lcampo] (\a:0.2) -- (\a:3.3);}
  \end{scope}
  \marcoL
  \node[cargapos] at (0,0) {$+$};
\end{tikzpicture}\\[2pt]{\small\scshape (a)\; carga puntual}
\end{minipage}\hfill
\begin{minipage}[t]{0.32\textwidth}\centering
\begin{tikzpicture}[x=0.8cm, y=0.8cm, baseline=(current bounding box.north)]
  \def\qs{-1}
  \begin{scope}
    \clip (-3,-3) rectangle (3,3);
    % equipotenciales (positivas: línea continua; negativas, por simetría: discontinua)
    \foreach \xs in {-1.25,-1.4,-1.65,-2.0,-2.6}{
      \equipC{\xs}{0}{0.05}{140}{-1}
      \draw[lequi] \ruta; \begin{scope}[yscale=-1]\draw[lequi] \ruta;\end{scope}
      \begin{scope}[xscale=-1]\draw[lequi, densely dashed] \ruta;\end{scope}
      \begin{scope}[xscale=-1,yscale=-1]\draw[lequi, densely dashed] \ruta;\end{scope}}
    \draw[lequi] (0,-3) -- (0,3);
    % líneas de campo (8 conos de igual flujo desde la carga positiva)
    \foreach \k in {1,...,8}{
      \pgfmathsetmacro{\psi}{acos(1-(2*\k-1)/8)}
      \pgfmathsetmacro{\sx}{-1+0.13*cos(\psi)}\pgfmathsetmacro{\sy}{0.13*sin(\psi)}
      \lineaC{\sx}{\sy}{0.05}{200}{1}
      \draw[lcampo] \ruta;
      \begin{scope}[yscale=-1]\draw[lcampo] \ruta;\end{scope}}
  \end{scope}
  \marcoL
  \node[cargapos] at (-1,0) {$+$};
  \node[carganeg] at (1,0) {$-$};
\end{tikzpicture}\\[2pt]{\small\scshape (b)\; dipolo}
\end{minipage}\hfill
\begin{minipage}[t]{0.32\textwidth}\centering
\begin{tikzpicture}[x=0.8cm, y=0.8cm, baseline=(current bounding box.north)]
  \def\qs{1}
  \begin{scope}
    \clip (-3,-3) rectangle (3,3);
    \foreach \xs in {-1.2,-1.3,-1.45,-1.8,-2.2,-2.7}{
      \equipC{\xs}{0}{0.05}{150}{-1}
      \foreach \sxx/\syy in {1/1,1/-1,-1/1,-1/-1}{\begin{scope}[xscale=\sxx,yscale=\syy]\draw[lequi] \ruta;\end{scope}}}
    \foreach \k in {1,...,8}{
      \pgfmathsetmacro{\psi}{acos(1-(2*\k-1)/8)}
      \pgfmathsetmacro{\sx}{-1+0.13*cos(\psi)}\pgfmathsetmacro{\sy}{0.13*sin(\psi)}
      \lineaC{\sx}{\sy}{0.05}{200}{1}
      \foreach \sxx/\syy in {1/1,1/-1,-1/1,-1/-1}{\begin{scope}[xscale=\sxx,yscale=\syy]\draw[lcampo] \ruta;\end{scope}}}
  \end{scope}
  \marcoL
  \node[cargapos] at (-1,0) {$+$};
  \node[cargapos] at (1,0) {$+$};
\end{tikzpicture}\\[2pt]{\small\scshape (c)\; dos cargas iguales}
\end{minipage}
}%
\caption{Líneas de campo (púrpura, con el sentido del campo) y superficies equipotenciales (azul, vistas en
un plano meridiano) de (a)~una carga puntual, (b)~un dipolo y (c)~dos cargas iguales. Las líneas se trazan con
igual flujo entre dos consecutivas, de modo que su densidad es proporcional a $\norm{\vE}$. Salen de las cargas
positivas y llegan a las negativas, nunca se cortan y son perpendiculares a las equipotenciales, hacia donde
el potencial disminuye. En (b) las equipotenciales de potencial negativo se trazan a trazos y la mediatriz es la
de potencial nulo; en (c) hay un punto de campo nulo en el punto medio.}
\label{fig:lineas}
\end{figure}
